\subsubsection*{1. Polar version}

\begin{description}
\item[(a)] \textbf{Rewrite the orbit equation using polar coordinate}
  \begin{align*}
    x &= R(\alpha)\cos\alpha \\
    y &= R(\alpha)\sin\alpha
  \end{align*}
  Thus, if we let
  \begin{align*}
    x^* &= R(\alpha)\cos\left(\alpha - \frac{\pi}{3}\right)
      = \frac{x}{2} + \frac{y\sqrt{3}}{2} \\
    y^* &= R(\alpha)\sin\left(\alpha - \frac{\pi}{3}\right)
      = -\frac{x\sqrt{3}}{2} + \frac{y}{2}
  \end{align*}
  \hfill(10 points)

  the equation of the orbit can be written as
  \begin{align*}
    9\left(R(\alpha)\cos\left(\alpha - \frac{\pi}{3}\right) - 4\right)^2
      + 25\left(R(\alpha)\sin\left(\alpha - \frac{\pi}{3}\right)\right)^2
      &= 225 \\
    9(x^* - 4)^2 + 25(y^*)^2 &= 225
  \end{align*}
  \hfill(10 points)

  \textbf{Determine semi major and semi minor axis of the orbit}

  Since
  \[
    9(x^* - 4)^2 + 25(y^*)^2 = 225
  \]
  is equivalent to
  \[
    \frac{(x^* - 4)^2}{25} + \frac{(y^*)^2}{9} = 1
  \]
  then, the semi major axis is 5, and \hfill(5 points) \\
  the semi minor axis is 3 \hfill(5 points)

\item[(b)] \textbf{Determine the zenith angle at perigee} \hfill(20 points)

  At the perigee, $\alpha - \frac{\pi}{3} = \pi$. Thus, the zenith angle at
  the perigee is
  \[
    \alpha - \frac{\pi}{2} = \pi + \frac{\pi}{3} - \frac{\pi}{2}
      = \frac{5\pi}{6}
  \]

% FIGURE NEEDED: rotated ellipse with the zenith angle at perigee marked at
% the origin (2015 long-question solutions, page 2/13).

\item[(c)] \textbf{Characterize the coordinates of the moon when it looks
  largest to the observer.}

  Let $P(x_0, y_0)$ be the point where it looks largest to the observer. Since
  the distance between the moon and the observer decreases as it moves from
  apogee towards perigee, then the moon looks largest to the observer at the
  planet when it closest to the planet while the whole of the moon still can % sic: "when it closest"
  be seen in full. Hence, $y_0 = r$.
  \begin{align*}
    y_0 &= r = R(\pi - \alpha_0\alpha_0)\sin(\pi - \alpha_0)
      = R_0\sin(\pi - \alpha_0) \\ % sic: "R(pi - alpha_0 alpha_0)" as printed
    x_0 &= R_0\cos(\pi - \alpha_0)
  \end{align*}
  \hfill(10 points)

% FIGURE NEEDED: sketch of the orbit arc near perigee with the moon of radius
% r at (x_0, r), the distance R_0 to the observer at the origin and the angle
% alpha_0 = theta_0/2 below the x-axis (2015 long-question solutions, page
% 3/13).

  \textbf{Determine the coordinates of the point.}

  Let the coordinate be $(x_0, r)$.
  \[
    9\left(\frac{x_0}{2} + \frac{\sqrt{3}r}{2} - 4\right)^2
      + 25\left(-\frac{\sqrt{3}x_0}{2} + \frac{r}{2}\right)^2 = 225
  \]
  Rewrite it as a quadratic equation in $x_0$.
  \begin{align*}
    9\left(x_0 + r\sqrt{3} - 8\right)^2
      + 25\left(-x_0\sqrt{3} + r\right)^2 &= 900 \\
    84(x_0)^2 + \left(-32\sqrt{3}r - 144\right)x_0
      + \left(52r^2 - 144\sqrt{3}r + 576 = 900\ \right) % sic: stray "= 900" inside the bracket as printed
  \end{align*}
  \hfill(20 points)

  Solving for $x$, we get
  \[
    x_0 = \frac{4r\sqrt{3} + 18}{21}
      \pm \frac{15\sqrt{-r^2 + 4r\sqrt{3} + 9}}{21}
  \]
  Then choose the smaller one
  \[
    x_0 = \frac{4r\sqrt{3} + 18}{21}
      - \frac{15\sqrt{-r^2 + 4r\sqrt{3} + 9}}{21}
  \]
  \hfill(10 points)

  \textbf{Find $\tan\frac{\theta}{2}$.} Hence
  \[
    \tan\frac{\theta_0}{2} = \tan\alpha_0 = \frac{r}{x_0}
      = \frac{21r}{4r\sqrt{3} + 18 - 15\sqrt{-r^2 + 4r\sqrt{3} + 9}}
  \]
  \hfill(10 points)
\end{description}

\subsubsection*{2. Cartesian version}

\begin{description}
\item[(a)] \textbf{Notice the standard version of the orbits}

  Simplify the equation by introducing new variables
  \begin{align*}
    x^* &= \frac{x}{2} + \frac{\sqrt{3}y}{2}
      = x\cos\frac{\pi}{3} + y\sin\frac{\pi}{3} \\
    y^* &= -\frac{\sqrt{3}x}{2} + \frac{y}{2}
      = x\sin\left(-\frac{\pi}{3}\right) + y\cos\frac{\pi}{3}
  \end{align*}
  \hfill(10 points)

  Write the orbit equation in terms new variables % sic: "in terms new variables"
  \[
    9(x^* - 4)^2 + 25(y^*) = 225 % sic: exponent missing on (y*) as printed
  \]
  \hfill(10 points)

% FIGURE NEEDED: plot of the original tilted ellipse through (x,y) and the
% rotated standard ellipse (dashed) through (x*,y*), axes graduated 0..8
% (2015 long-question solutions, page 4/13).

  \textbf{Determine semi major and semi minor axis of the orbit}

  Since
  \[
    9(x^* - 4)^2 + 25(y^*)^2 = 225
  \]
  is equivalent to
  \[
    \frac{(x^* - 4)^2}{25} + \frac{(y^*)^2}{9} = % sic: right-hand side left blank as printed
  \]
  then, the semi major axis is 5, and \hfill(5 points) \\
  the semi minor axis is 3 \hfill(5 points)

\item[(b)] \textbf{Determine the zenith angle at perigee} \hfill(20 points)

  Since the original ellipse may be obtained from standard ellipse
  \[
    \frac{(x^* - 4)^2}{25} + \frac{(y^*)^2}{9} = 1
  \]
  by $\pi/3$ counterclockwise rotation respect to the origin, the the zenith % sic: "the the"
  angle at perigee is equal to
  \[
    \pi + \frac{\pi}{3} - \frac{\pi}{2} = \frac{5\pi}{6}
  \]

\item[(c)] \textbf{Characterize the coordinates of the moon when it looks
  largest to the observer.}

  Let $P(x_0, y_0)$ be the point where it looks largest to the observer. Since
  the distance between the moon and the observer decreases as it moves from
  apogee towards perigee, then the moon looks largest to the observer at the
  planet when it closest to the planet while the whole of the moon still can
  be seen in full. Hence, $y_0 = r$. \hfill(10 points)

  \textbf{Find $x_0$.}
  \[
    9\left(\frac{x_0}{2} + \frac{\sqrt{3}r}{2} - 4\right)^2
      + 25\left(-\frac{\sqrt{3}x_0}{2} + \frac{r}{2}\right)^2 = 225
  \]
  Rewrite it as a quadratic equation in $x$.
  \begin{align*}
    9\left(x_0 + r\sqrt{3} - 8\right)^2
      + 25\left(-x_0\sqrt{3} + r\right)^2 &= 900 \\
    \left(9x_0^2 + \left(18\sqrt{3}r - 144\right)x_0 + 27r^2
      - 144\sqrt{3}r + 576\right)
      + \left(75x_0^2 - 50\sqrt{3}rx_0 + 25r^2\right) &= 900 \\
    84(x_0)^2 + \left(-32\sqrt{3}r - 144\right)x_0
      + \left(52r^2 - 144\sqrt{3}r + 576\ \right) &= 900
  \end{align*}
  Solving it for $x$ we get
  \[
    x_0 = \frac{4r\sqrt{3} + 18}{21}
      \pm \frac{15\sqrt{-r^2 + 4r\sqrt{3} + 9}}{21}
  \]
  \hfill(20 points)

  Then choose the smaller one
  \[
    x_0 = \frac{4r\sqrt{3} + 18}{21}
      - \frac{15\sqrt{-r^2 + 4r\sqrt{3} + 9}}{21}
  \]
  \hfill(10 points)

  \textbf{Find $\tan\frac{\theta}{2}$.} Hence
  \[
    \tan\frac{\theta}{2} = \tan\alpha_0 = \frac{r}{x_0}
      = \frac{21r}{4r\sqrt{3} + 18 - 15\sqrt{-r^2 + 4r\sqrt{3} + 9}}
  \]
  \hfill(10 points)
\end{description}

\subsubsection*{3. Cartesian version} % sic: both the second and third versions are titled "Cartesian version"

\begin{description}
\item[(a)] \textbf{Notice the standard version of the orbits}

  Simplify the equation by introducing new variables
  \begin{align*}
    x^* &= \frac{x}{2} + \frac{\sqrt{3}y}{2}
      = x\cos\frac{\pi}{3} + y\sin\frac{\pi}{3} \\
    y^* &= -\frac{\sqrt{3}x}{2} + \frac{y}{2}
      = x\sin\left(-\frac{\pi}{3}\right) + y\cos\frac{\pi}{3}
  \end{align*}
  \hfill(10 points)

  Write the orbit equation in terms new variables
  \[
    9(x^* - 4)^2 + 25(y^*) = 225
  \]
  \hfill(10 points)

  \textbf{Determine semi major and semi minor axis of the orbit}

  Since
  \[
    9(x^* - 4)^2 + 25(y^*)^2 = 225
  \]
  is equivalent to
  \[
    \frac{(x^* - 4)^2}{25} + \frac{(y^*)^2}{9} =
  \]
  The semi major axis is 5 \hfill(5 points) \\
  The semi major axis is 3 \hfill(5 points) % sic: "semi major" twice as printed

\item[(b)] \textbf{Notice the standard version of the orbits} \hfill(20 points)
  % sic: heading repeated from (a) as printed

  Since the original ellipse may be obtained from standard ellipse
  \[
    \frac{(x^* - 4)^2}{25} + \frac{(y^*)^2}{9} = 1
  \]
  by counterclockwise $\pi/3$ rotation respect to the origin, the zenith angle
  at perigee is
  \[
    \pi + \frac{\pi}{3} - \frac{\pi}{2} = \frac{5\pi}{6}
  \]

\item[(c)] \textbf{Identify and characterized the coordinates when the planet
  looks largest.} % sic: "characterized"

  Since the distance between the moon and the observer decreases as it moves
  from apogee towards perigee, then the moon looks largest to the observer at
  the planet when it closest to the planet while the whole of the moon still
  can be seen in full. Hence, it is at necessary that $y = r$. Let the % sic: "it is at necessary"
  coordinates be $(x_0, r)$. \hfill(10 points)

% FIGURE NEEDED: same near-perigee sketch as in the polar version, now with
% the moon at (x,r), distance d_theta to the observer and angle alpha =
% theta/2 (2015 long-question solutions, page 6/13).

  \textbf{Obtain expression for calculating distance from a point on an
  ellipse to the foci to set up an equation}

  Consider standard ellipse
  \[
    \frac{(x^* - 4)^2}{5^2} + \frac{(y^*)^2}{3^2} = 1
  \]
  Choose any point $(x^*, y^*)$ on the ellipse. Let $d_1$ and $d_2$ be the
  distances from any point $(x^*, y^*)$ on the ellipse to the foci $(0{,}0)$
  and $d(8{,}0)$; $d_1^2 = (x^*)^2 + (y^*)^2$ and
  $d_2^2 = (x^* - 8)^2 + (y^*)^2$

  Thus,
  \[
    d_1^2 = (x^*)^2 + (y^*)^2
      = (x^*)^2 + \left(9 - \frac{9(x^* - 4)^2}{5^2}\right)
      = \frac{(16(x^*)^2 + 72x^* + 81)}{25}
  \]
  Therefore
  \[
    d_1^2 = \frac{(16(x^*)^2 + 72x^* + 81)}{25} = x^2 + r^2
  \]
  \hfill(10 points)

  where $x^* = \frac{x}{2} + \frac{r\sqrt{3}}{2}$. Therefore, we can obtain
  the value of $x$ by solving the above equation for $x$.

  \textbf{Solve equation
  $\frac{(16(x^*)^2 + 72x^* + 81)}{25} = x^2 + r^2$ to obtain value of $x$.}

  Substituting $x^* = \frac{x}{2} + \frac{r\sqrt{3}}{2}$ to the equation, we
  have
  \[
    25d_1^2 = 25x^2 + 25r^2
      = 16\left(\frac{x}{2} + \frac{r\sqrt{3}}{2}\right)^2
      + 72\left(\frac{x}{2} + \frac{r\sqrt{3}}{2}\right) + 81 = % sic: trailing "=" as printed
  \]
  or
  \begin{align*}
    25x^2 + 25r^2 &= 4x^2 + \left(8r\sqrt{3} + 36\right)x
      + \left(12r^2 + 36\sqrt{3}r + 81\right) \\
    21x^2 - \left(8r\sqrt{3} + 36\right)x
      + \left(13r^2 - 36r\sqrt{3} - 81\right) &= 0
  \end{align*}
  Then use quadratic formula to obtain $x$ in term of $r$. Choose smaller
  value of $x$ to get larger value of
  $\tan\frac{\theta}{2} = \frac{r}{x}$. The smaller one is
  \[
    x = \frac{4r\sqrt{3} + 18}{21}
      - \frac{15\sqrt{-r^2 + 4r\sqrt{3} + 9}}{21}
  \]
  \hfill(20 points)

  \textbf{Find the expression of $\tan\frac{\theta}{2}$.}

  Hence
  \[
    \tan\frac{\theta}{2} = \frac{r}{x}
      = \frac{21r}{4r\sqrt{3} + 18 - 15\sqrt{-r^2 + 4r\sqrt{3} + 9}}
  \]
  \hfill(10 points)
\end{description}
